\section{Consciousness (\textit{Citta})}

Welcome to the third talk of this Practical Abhidhamma Course.\footnote{More details in Chapter 1 of “A Comprehensive Manual of Abhidhamma” (see Footnote 2 for link).} This talk will describe consciousness, the first of the four Ultimate Realities (\textit{citta}, \textit{cetasika}, \textit{rūpa} and \textit{Nibbāna}). You should have the \textit{Satipaṭṭhāna} Sutta printout and Handout 2 in front of you. 

Handout 2 lists Thought Moments in which consciousness arises, numbered from \textbf{1}--\textbf{89}. We will refer back to it many times during other talks as well. Consider it to be a map of the mind. I hope that with this map, you don’t get lost!

\subsection*{Definition of Consciousness}

Let’s start with a definition of consciousness. I downloaded the Wikipedia entry on “consciousness” and it filled 18 pages.\footnote{\url{http://en.wikipedia.org/wiki/Consciousness}} Obviously, “consciousness” is not a simple thing to define. I do, however, like the first sentence from the Wikipedia entry: “Consciousness is the quality or state of awareness, or, of being aware of an external object or something within oneself.”

There are two things to take away from this definition. First, consciousness and awareness are synonyms. Second, consciousness takes either an external object or an internal object. Examples of external objects include \textbf{Visible-form}, \textbf{Sound}, \textbf{Odour}, \textbf{Taste} and tactile object. When the mind is aware of itself, it takes an internal object such as \textbf{Delusion}, \textbf{Attachment}, \textbf{Faith} and \textbf{Compassion}.

So what is the Buddhist definition of “consciousness?” The Commentary defines consciousness as having three roles. First, consciousness is that which is aware of an object. There is no Self that is aware, it is consciousness that is aware. Second, consciousness allows the associated Mental Factors to be aware of an object. And third, consciousness is an activity, a process of being aware of an object.

\subsection*{Definition of Thought Moment}

Consciousness never arises alone; it always arises as part of a Thought Moment. This concept of a Thought Moment comes from the Abhidhamma, not from the Suttas.\footnote{Though the following beautiful poem from the Mahāniddesa describes a mind that is constantly changing, Nm 2.4: \url{http://www.accesstoinsight.org/tipitaka/kn/nm/nm.2.04.olen.html}} A Thought Moment is consciousness together with its associated Mental Factors.

Within a Thought Moment consciousness, and each of its associated Mental Factors, work as an inseparable team, each with their own role to play. Consciousness has the role of awareness of the object, and each Mental Factor has its own role.

As I mentioned during the first talk, according to the Commentaries, what is conventionally called the mind is actually a sequence of Thought moments. Each Thought Moment arises based on conditions, performs its function and then falls away again.

\pagebreak

\subsection*{Brief descriptions of each group in Handout 2}

Please look at Handout 2. In this “map of the mind,” Thought Moments are grouped together in rows and columns. I will first give a brief description of each group and then do an in-depth analysis of some of the groups. The in-depth analysis of the remaining groups will be done in later talks. We will refer to Handout 2 many times during these talks. Let’s start with the brief description of each group.

\begin{figure}[h]
\centering
\includegraphics[width=1\linewidth]{./Diagrams/Handout_2}
\caption{The Danger Zone, the Faultless Zone and the Sensing Zone are the most important areas of Handout 2.}
\label{fig:Handout_2}
\end{figure}

In the first group are Thought Moments \textbf{1}--\textbf{12}. These Thought Moments create new, unwholesome\footnote{The Pāḷi term for unwholesome is \textit{akusala}, the opposite of \textit{kusala}; the prefix “\textit{a}” signifies the “active opposite of,” not merely “absence of.” The Commentary defines \textit{kusala} as “healthy,” “faultless” and “producing happy results.” Therefore, \textit{akusala} means “unhealthy,” “faulty,” and “producing unhappy results.”} kamma. I will call this the “Danger Zone.” \footnote{Abhidhamma texts do not use the term, “Danger Zone.”} The Danger Zone is engulfed in a dense fog, so it is impossible for the mind to see where it is, or where it is going. The dense fog is \textbf{Delusion}. In part of the Danger Zone there is sticky quicksand that traps the mind, symbolizing \textbf{Attachment}-rooted Thought Moments. Another part of the Danger Zone is burning hot and painful for the mind; this represents \textbf{Aversion}-rooted Thought Moments. Unfortunately, the mind spends almost all of its time in the Danger Zone and therefore treats the Danger Zone as its home. The fog here is so thick that the mind does not even realize that this is a Danger Zone.

In the second group are Thought Moments \textbf{13}--\textbf{30}. These Thought Moments are involved in the sensing and processing of sense data and ideas. I will call this the “Sensing Zone.” \footnote{Abhidhamma texts do not use the term, “Sensing Zone.”}

In the third group are Thought Moments \textbf{31}--\textbf{38}. These create new, wholesome kamma. I will call this the “Faultless Zone.” \footnote{The Abhidhammattha Sangaha classifies Thought Moments \textbf{31}--\textbf{54} as being “beautiful” (\textit{sobhana}). Abhidhamma texts do not use the term, “Faultless Zone,” but “faultless” is used as a translation of “\textit{kusala}.”} In the Faultless Zone there is no \textbf{Attachment}, only \textbf{Non-attachment}. There is no \textbf{Aversion}, only \textbf{Non-aversion}. In some parts of the Faultless Zone, details of everything can be clearly seen, and these parts are the Thought Moments associated with \textbf{Understanding}.

In a Sutta called “Two Sorts of Thinking,” the Buddha stressed the importance of knowing if the mind was in the Danger Zone or in the Faultless Zone.\footnote{MN 19: \url{http://www.accesstoinsight.org/tipitaka/mn/mn.019.than.html}} If the mind is in the Danger Zone, the Buddha tells us to reflect on the disadvantages of such thinking. If the mind is in the Faultless Zone, just be passively aware. In this Sutta, the Buddha also said that whatever you keep thinking about will become the inclination of your awareness. In other words, unwholesome thinking accumulates into an unwholesome perspective and unwholesome habits.

\begin{figure}[h]
\centering
\includegraphics[width=1.0\linewidth]{./Diagrams/Life_Cont}
\caption{Many Life-continuum Thought Moments arise between each Sensing/Thinking Process and during dreamless sleep.}
\label{fig:Life_Cont}
\end{figure}

In the fourth group are Thought Moments \textbf{39}--\textbf{46}. The main function of these Thought Moments is “Life-continuum.” \footnote{The Pāḷi term for “Life-continuum” is \textit{bhavaṅga}; literally factor (\textit{aṅga}) of existence (\textit{bhava}). The only reference to \textit{bhavaṅga} in the Abhidhamma \textit{Piṭaka} is in “Conditional Relations” (\textit{Paṭṭhāna}), Volume I page 159 and elsewhere, where it is described as preceding five-sense-door adverting (Thought Moment \textbf{28}).} “Life-continuum” Thought Moments arise when the mind is not sensing, not thinking and not in a jhāna\footnote{\url{http://en.wikipedia.org/wiki/Dhyana_in_Buddhism}} meditative state. “Life-continuum” Thought Moments arise during dreamless sleep, and between all the sensing and thinking that goes on while we are awake.

In the fifth group are Thought Moments \textbf{47}--\textbf{54}. These arise only in an Arahat. Thought Moments \textbf{47}--\textbf{54} correspond to Thought Moments \textbf{31}--\textbf{38}, except that Thought Moments \textbf{47}--\textbf{54} do not create new kamma as Arahats\footnote{\url{http://en.wikipedia.org/wiki/Arhat}} do not create any new kamma.

In the sixth group are Thought Moments \textbf{55}--\textbf{59}. These arise during the five “Fine Material” jhāna and create new, wholesome kamma.

In the seventh group are Thought Moments \textbf{60}--\textbf{64}. These do not arise in humans; they are the “Life-continuum” Thought Moments for beings in higher Realms of Existence.

In the eighth group are Thought Moments \textbf{65}--\textbf{69}. These arise only in an Arahat,when an Arahat experiences a Fine Material jhāna.

In the ninth group are Thought Moments \textbf{70}--\textbf{73}. These arise during the four “Immaterial” jhāna and create new, wholesome kamma.

In the tenth group are Thought Moments \textbf{74}--\textbf{77}. These do not arise in humans; they are the “Life-continuum” Thought Moments for beings in even higher Realms of Existence.

In the eleventh group are Thought Moments \textbf{78}--\textbf{81}. These arise only when an Arahat experiences an Immaterial jhāna.

In the twelfth group are Thought Moments \textbf{82}--\textbf{85}. These “path” Thought Moments arise at the moment of transition to each degree of sainthood. In the thirteenth group Thought Moments \textbf{86}--\textbf{89}. These are the kammic result of Thought Moments \textbf{82}--\textbf{85}. The Pāḷi name for the twelfth and thirteenth groups is translated as “Supramundane,” literally means “transcending the world.” In other words, they are off the map. These are the only Thought Moments that can take \textit{Nibbāna} as object and are reserved for saints and for the Buddha.

\subsection*{Using the Handout}

Before we go into the details of Thought Moments, let’s take a step back and consider the question, “How can I use this handout?”

Following the advice of the Buddha, to develop spiritually we need to look at the mind. If the mind is agitated, we need first to steady, settle, unify and compose the mind; perhaps by focusing on a neutral object such as the breath. When we look inwards with a calm mind, the most obvious Thought Moments are the Danger Zone (Thought Moments \textbf{1}--\textbf{12}) and the Faultless Zone (Thought Moments \textbf{31}--\textbf{38}).

If unwholesome Thought Moments are found when looking at the mind, don’t be frustrated. Frustration is another unwholesome Thought Moment. If unwholesome Thought Moments are found, be mindful of their characteristics; “Ah, this is the grasping of \textbf{Attachment}.” or “Ah, the mind is experiencing \textbf{Aversion}.” Mindfulness converts the unwholesome Thought Moment into a wholesome Thought Moment with \textbf{Understanding}.

If wholesome Thought Moments are found when looking at the mind, don’t be attached to them; craving wholesome Thought Moments brings craving, it is not a wholesome Thought Moment. If wholesome Thought Moments are found, be mindful of their characteristics; “Ah, the mind is calm, not floating away.”

Here is a \textbf{\textit{RADICAL}} approach based on a modern author.\footnote{Sayālay Susīlā: \url{http://www.sayalaysusila.net/files/Moment-to-Moment-Practice.pdf}} It is a traditional approach, but the word \textbf{\textit{RADICAL}} is used as an acronym. \textbf{\textit{R}} stands for Recognize: identify the current Thought Moment or physical sensation being experienced; \textbf{Sound}, \textbf{Hardness}, \textbf{Heat}. These are examples of physical sensations. \textbf{\textit{A}} stands for Accept: accept what is experienced just as it is; resisting the unpleasant is \textbf{Aversion}, and clinging to the pleasant is \textbf{Attachment}. \textbf{\textit{D}} stands for Depersonalize: become a stable observer, separate and unbiased; what is experienced is not “happening to me,” it is not “myself,” it is not “mine,” it simply is something to be observed. \textbf{\textit{I}} stands for Investigate. This is the point at which \textbf{Understanding} starts to play an important role. Look closely; what is experienced has characteristics. Investigate these characteristics until they are understood. \textbf{\textit{CA}} stands for Contemplate \textit{Anicca}. Notice that what is experienced is impermanent. \textbf{\textit{CA}} can also stand for Contemplate \textit{Anattā}. Notice that what is experienced is a natural process, not Self. \textbf{\textit{L}} stands for Let go; there is always a new experience waiting. So the \textbf{\textit{RADICAL}} approach is \textbf{\textit{R}}ecognize, \textbf{\textit{A}}ccept, \textbf{\textit{D}}epersonalize, \textbf{\textit{I}}nvestigate, \textbf{\textit{C}}ontemplate \textbf{\textit{A}}nicca or \textbf{\textit{C}}ontemplate \textbf{\textit{A}}nattā and \textbf{\textit{L}}et go.

\pagebreak

\subsection*{Thought Moments 1--12 (Danger Zone)}

\begin{figure}[h]
\centering
\includegraphics[width=0.45\linewidth]{./Diagrams/Danger}
\caption{Thought Moments \textbf{1}--\textbf{12}: Danger Zone}
\label{fig:Danger}
\end{figure}

Let’s look at the Danger Zone, Thought Moments \textbf{1}--\textbf{12}, in more detail. This group is subdivided into “\textbf{Attachment}-rooted,” “\textbf{Aversion}-rooted” and “\textbf{Delusion}-rooted.” According to the Suttas, there are three unwholesome roots: \textbf{Attachment} (\textit{lobha}), \textbf{Aversion} (\textit{dosa}) and \textbf{Delusion} (\textit{moha}),\footnote{Iti 3.50: \url{http://www.accesstoinsight.org/tipitaka/kn/iti/iti.3.050-099.than.html\#iti-050}} and these Mental Factors arise in certain Thought Moments. Thought Moments \textbf{1}--\textbf{12} include either one or two unwholesome roots. Just as roots provide a foundation, support, strength and nourishment to a tree, the roots of a Thought Moment provide a foundation, support, strength and nourishment to a Thought Moment.

A group of laypeople asked the Buddha how to evaluate different teachings.\footnote{AN 3.65: \url{http://www.accesstoinsight.org/tipitaka/an/an03/an03.065.than.html}} The Buddha first advised them not to judge teachings based on tradition, conjecture or the charisma of the teacher.\footnote{At the time of the Buddha, there were three approaches to knowledge: oral tradition (from the ancient Vedas: \url{http://en.wikipedia.org/wiki/Vedas}), logical reasoning (the basis of the recently-developed Upaniṣads: \url{http://en.wikipedia.org/wiki/Upanishads}, teachings of religious philosophy) and direct intuition of a teacher. In MN 100 (\url{http://metta.lk/tipitaka/2Sutta-Pitaka/2Majjhima-Nikaya/Majjhima2/100-sangarava-e1.html}), the Buddha put himself in the third category. See \url{http://www.ahandfulofleaves.org/documents/Early Buddhist Theory of Knowledge_Jayatilleke.pdf}} The Buddha then asked if \textbf{Attachment}, \textbf{Aversion} and \textbf{Delusion} were beneficial or unbeneficial, and if \textbf{Non-attachment}, \textbf{Non-aversion} and \textbf{Understanding} were beneficial or unbeneficial. The laypeople said that \textbf{Attachment}, \textbf{Aversion} and \textbf{Delusion} were unbeneficial and \textbf{Non-attachment}, \textbf{Non-aversion} and \textbf{Understanding} were beneficial. The Buddha then advised that teachings were good if they led to a decrease in \textbf{Attachment}, \textbf{Aversion} and \textbf{Delusion} and led to an increase in \textbf{Non-attachment}, \textbf{Non-aversion} and \textbf{Understanding}. This simple guideline can be applied when considering the motivation behind our own actions.

\subsubsection*{\textbf{Attachment}-rooted}

Thought Moments \textbf{1}--\textbf{8} are \textbf{Attachment}-rooted. Earlier, I referred to these Thought Moments as the sticky quicksand that traps the mind in the foggy Danger Zone. Sticky quicksand refers to the root of \textbf{Attachment} which is the primary characteristic of these Thought Moments. These eight Thought Moments also have a root of \textbf{Delusion}; this is the thick fog. They are in the Danger Zone because they create unwholesome kamma.

Reminds me of a joke... the young novice monk asks his teacher, “Master, is it allowable for a monk to use email?” The teacher replies, “Sure, as long as there are no attachments.” \footnote{\textbf{Attachment} to email is a significant distraction during a meditation retreat. Yogis may want to turn off data (to avoid email and instant messages) and tell people to send an SMS in case of emergency.}

\textbf{Attachment} has many grades ranging from simply enjoying a morning coffee to intense lustful passion. A subtle form of \textbf{Attachment} is \textbf{Attachment} to sense objects; the desire to see things, to hear things, etc. There is a story in the Suttas\footnote{SN 47.7: \url{http://www.accesstoinsight.org/tipitaka/sn/sn47/sn47.007.than.html}} about sense-desires.\footnote{Suttas warning of the dangers of sense-desires are directed at monks, who have already overcome most of the attachments faced by laypeople.} A hunter puts tar on a branch. The foolish monkey is curious and touches the tar with one hand. The hand gets stuck, so the monkey uses his other hand, his two feet and finally his mouth to try to get free. The monkey is then trapped in five ways, corresponding to the five senses-desires. I like the modern version of this story\footnote{See \url{http://en.wiktionary.org/wiki/monkey_trap}} where the hunter attaches a hollowed-out coconut to a tree. Food is put inside the coconut to attract the monkey. The monkey grabs the food but with a clenched fist full of food, the monkey is unable to extract its hand from the coconut so he is trapped. I like this version because to free itself, the monkey just has to overcome its instinctive \textbf{Attachment} to the food and open its hand. Of course, overcoming instinctive \textbf{Attachment} is easier said than done for the monkey, and for the monkey’s distant cousins, we humans!

The eight Thought Moments in this group include every combination of “accompanied by pleasant or indifferent \textbf{Feeling},” “associated or not associated with \textbf{Wrong view}” and “unprompted or prompted.”

The first differentiating factor is accompanied by either pleasant or indifferent \textbf{Feeling}. Dictionary.com\footnote{\url{http://dictionary.reference.com/browse/feeling}} has 14 definitions for “feeling.” In the Abhidhamma, the word “\textbf{Feeling}” has one definition; the most basic experience of pleasant, unpleasant or neutral. When we see the word “\textbf{Feeling}” in the Suttas or Abhidhamma, it is important that we put aside all other meanings associated with the English word “feeling,” and think of “\textbf{Feeling}” as describing only a pleasant, unpleasant or neutral experience. Later, we will discuss the contemplation of \textbf{Feeling} in the \textit{Satipaṭṭhāna} Sutta.

The second differentiating factor, either “associated with \textbf{Wrong view}” or “not associated with \textbf{Wrong view},” is extremely important, so let’s discuss \textbf{Attachment} to \textbf{Wrong view} in a bit more detail. \textbf{Attachment} to \textbf{Wrong view} means clinging to an opinion, a bias or a judgement. \textbf{Attachment} to \textbf{Wrong view} is one’s outlook, perspective, paradigm or belief.

Almost 500 years ago, Copernicus\footnote{\url{http://en.wikipedia.org/wiki/Nicolaus_Copernicus}} wrote that observations and calculations showed that the earth orbited the sun; that the earth is not the centre of the universe. Initially, the idea was ignored. Sixty years after Copernicus’ death, only 15 astronomers in all of Europe agreed. The idea was ignored because it challenged the prevailing “common sense” and church doctrine.

In my first talk, I shared a scientific experiment showing that a “Self who decides” is an illusion created after a decision has already been made; there is no Self at the centre of the universe. I suspect that after an initial discomfort with the results of the experiment, your mind ignored results of the experiment because the ego-centric view of the universe is “common sense.” “Ignored” is a verb; the related noun is “ignorance” or mental blindness. Refusing to see things as they truly are is clinging to ignorance or \textbf{Attachment} to \textbf{Wrong view}.

But we don’t need a fancy experiment to show the illusion of a controlling Self. Consider the last time that confusion arose in the mind. Was there a conscious decision at that time? Was there a Self who said, “I think that I will choose to be confused now?” Of course not! Confusion arose naturally. After confusion arose, Self-identification took place, “I am confused.”

\begin{figure}[h]
\centering
\includegraphics[width=1\linewidth]{./Diagrams/Latent}
\caption{\textit{\small Process by which latent defilements are reinforced.}}
\label{fig:Latent}
\end{figure}

Imagine a man with a fear of snakes.\footnote{This famous analogy is used by Candrakirti in his Commentary on Aryadeva’s “Catuḥśataka”.} Most of the time, this fear has no impact on his life. The fear is latent, not arisen, and he may not even be aware of this fear. The man is walking down a dimly-lit path at night and there is a coil of rope ahead. His latent fear of snakes causes him to misperceive the coil as a snake. This distorted perception, combined with his latent fear of snakes, causes him to think about snakes. This distorted thinking, combined with his latent fear of snakes, convinces the man that there is a snake ahead. This \textbf{Wrong view} causes the fear of snakes to arise. It is no longer a latent fear. This actual fear reinforces his latent tendency to be afraid of snakes.\footnote{Distortion of \textbf{Perception} is \textit{saññāvipallāsa}, distortion of thought is \textit{cittavipallāsa} and distortion of view is \textit{diṭṭhivipallāsa}. Others have translated \textit{vipallāsa} as “inversions” or “perversions;” \textit{Vipallāsa} is from \textit{vi} + \textit{pari} + \textit{āsa} = “turned upside down.”} In other words, actual fear comes from latent fear plus distorted view, distorted view comes from latent fear plus distorted thought, distorted thought comes from latent fear plus distorted perception, and distorted perception comes from latent fear.

Imagine a woman who is attached to the idea of spiritual progress. This is a latent \textbf{Attachment}, waiting for a trigger to  arise. While sitting on her cushion one day, some mundane experience arises in the woman’s mind. Her \textbf{Attachment} to spiritual progress causes her to misperceive the mundane experience as an insight. Driven by this distorted perception, and supported by her \textbf{Attachment} to the idea of spiritual progress, she thinks she has experienced an insight. This distorted thought, supported by her \textbf{Attachment} to the idea of spiritual progress causes the woman to be convinced that she has experienced an insight. This \textbf{Wrong view} causes \textbf{Attachment} to arise. It is no longer a latent \textbf{Attachment}. This actual \textbf{Attachment} reinforces her latent tendency to be attached to the idea of spiritual progress.

\textbf{Delusion} is mental blindness, while \textbf{Wrong view} involves conviction that something that is wrong is correct. As you can see from the examples, \textbf{Delusion} triggers distorted perception, \textbf{Delusion} and distorted perception trigger distorted thought, then \textbf{Delusion} and distorted thought trigger \textbf{Attachment} to \textbf{Wrong view}. When \textbf{Attachment} to \textbf{Wrong view} arises, the underlying \textbf{Delusion} is reinforced.

\pagebreak

\begin{figure}[h]
\centering
\includegraphics[width=0.7\linewidth]{./Diagrams/Cats}
\caption{\textbf{Delusion} creates a distortion. \textbf{Attachment} to \textbf{Wrong view} sees something that does not really exist (for example, the ego sees a huge “Self”).}
\label{fig:Cats}
\end{figure}

\begin{figure}[h]
\centering
\includegraphics[width=1\linewidth]{./Diagrams/Perception}
\caption{Latent tendencies influence how an image is perceived; when looking at the first image, do you perceive a vase or two faces? Psychologists use the Rorschach inkblot test (\url{http://en.wikipedia.org/wiki/Rorschach_test}), such as the second image, to gain insights into a person’s latent tendencies.}
\label{fig:Perception}
\end{figure}

The Buddha listed four types of \textbf{Attachment} to \textbf{Wrong view} that are harmful to spiritual progress.\footnote{AN 4.49: \url{http://www.accesstoinsight.org/tipitaka/an/an04/an04.049.than.html}} First, taking what is impermanent as permanent is \textbf{Wrong view}. Second, taking what is unsatisfactory as pleasant is \textbf{Wrong view}. Third, taking what is impure to be pure is \textbf{Wrong view}. Fourth, taking what is non-self as Self is \textbf{Wrong view}. The Buddha frequently highlighted the danger of taking what is non-self as Self.\footnote{According to Visuddhimagga XXII.68 (see footnote 2) the Sotāpanna uproots all four types of \textit{diṭṭhivipallāsa} (this implies that no speech or action will take place because of the \textit{vipallāsa}). \textit{Saññāvipallāsa} and \textit{cittavipallāsa} related to \textit{anicca} is uprooted by the Sotāpanna, \textit{saññāvipallāsa} and \textit{cittavipallāsa} related to \textit{dukkha} is uprooted by the Arahat, \textit{saññāvipallāsa} and \textit{cittavipallāsa} related to \textit{asubha} is uprooted by the Anāgāmī, \textit{saññāvipallāsa} and \textit{cittavipallāsa} related to \textit{anattā} is uprooted by the Sotāpanna.} He explained how the ego of the untrained mind twists and distorts whatever is perceived, to place a Self at the centre by imagining a relationship with what is experienced.\footnote{MN 1: \url{http://www.accesstoinsight.org/tipitaka/mn/mn.001.than.html}} The experience of pain becomes “my pain,” “I am in pain” or “it is painful to me.” All of this happens in the blink of an eye.

The Buddha’s teachings, our own experience, and scientific experiments all tell us that “Self is an illusion.” In spite of all this evidence, we stubbornly cling to the \textbf{Wrong view} of Self. I hope that reflecting on the content of this Practical Abhidhamma Course, and putting it into practice, helps to loosen the grip that ego has on your mind.

The third differentiating factor is unprompted or prompted. An unprompted Thought Moment arises spontaneously whereas a prompted Thought Moment has to be instigated or induced. If somebody convinces me to do something, if I have to convince myself to do it or if I do it reluctantly, then the Thought Moment behind the action is prompted, instigated or induced. On the other hand, if I do something spontaneously, without anybody trying to convince me, then the Thought Moment behind the action is unprompted. Prompting is not based on the object, prompting is related to influence by another person or oneself.

Consider a boy spontaneously stealing an apple, with joy, convinced that there is nothing wrong with stealing. The boy’s Thought Moment is unprompted, associated with \textbf{Wrong view} and accompanied by pleasant \textbf{Feeling}. It is Thought Moment \textbf{1}.

Imagine that, knowing lying is wrong, I reluctantly compliment someone to make them happy. My Thought Moment is prompted, not associated with \textbf{Wrong view}, and accompanied by indifferent \textbf{Feeling}. It is Thought Moment \textbf{8}.

The \textbf{Attachment}-rooted Thought Moments are subdivided according to prompting, \textbf{Wrong view} and \textbf{Feeling}. This subdivision may have been done to reflect the relative weightiness of the associated kamma. A Thought Moment that is unprompted or spontaneous will create weightier kamma than a prompted Thought Moment that depends on somebody else, or depends on self-reflection. A Thought Moment associated with \textbf{Wrong view} will create weightier kamma than a Thought Moment not associated with \textbf{Wrong view}.\footnote{For example, Milindapañha III,7,8 explains that a person who picks up a hot object unknowingly will suffer greater burns than a person who picks up a hot object knowingly.} A Thought Moment accompanied by pleasant \textbf{Feeling} will create weightier kamma than a Thought Moment accompanied by indifferent \textbf{Feeling}. Spontaneity, association with \textbf{Wrong view} and pleasant \textbf{Feeling} all contribute to increase the \textbf{Volition} or intention of a Thought Moment, and weightiness of kamma depends on the level of \textbf{Volition} or intention.

I am not saying that stealing the apple resulted in weightier kamma than giving a false compliment. What I am saying, is that stealing the apple spontaneously created weightier kamma than if the boy had to be convinced by his friend to steal it. The weightiness of kamma depends on the strength of the underlying \textbf{Volition}, so it is difficult to compare different actions.

This weightiness of kamma has a parallel in how our legal system works. Is a judge not more likely to give a stiffer punishment if a defendant is happy they committed a crime or think that they did nothing wrong? Is a judge not more likely to be more lenient if a defendant was forced to commit the crime or committed the crime reluctantly?

\subsubsection*{\textbf{Aversion}-rooted}

Now let’s look at Thought Moments rooted in \textbf{Aversion}. Thought Moment \textbf{9} is unprompted and Thought Moment \textbf{10} is prompted. Earlier, I referred to these Thought Moments as burning hot and painful because they are accompanied by unpleasant \textbf{Feeling}. While the eight \textbf{Attachment}-rooted Thought Moments are attracted to the object, these two Thought Moments do not like the object or are not satisfied with the object.

Just as \textbf{Attachment} has many grades, from simply enjoying a morning coffee to intense lustful passion, \textbf{Aversion} also has many grades. \textbf{Aversion} could be as subtle as not accepting the object as it is, or as intense as blinding hatred. When there is fear, there is \textbf{Aversion} towards a future situation. When there is \textbf{Remorse} or regret, there is \textbf{Aversion} towards a past situation. When there is boredom, there is \textbf{Aversion} towards the current situation. When there is \textbf{Envy}, there is \textbf{Aversion} because I do not like it that somebody has something that I do not. Hatred, anger, fear, \textbf{Remorse}, boredom and \textbf{Envy} are all accompanied by unpleasant \textbf{Feeling}.

Which Thought Moment arises at the moment a hunter kills for sport? Before the killing, the hunter was looking for the animal and the Thought Moments were rooted in craving or \textbf{Attachment}. At the moment of killing, the hunter has \textbf{Aversion} to the fact that the animal is alive. At this moment, the hunter does not want the animal to continue living, so the Thought Moment is rooted in \textbf{Aversion}. Since the killing is premeditated, the Thought Moment is prompted. This is Thought Moment \textbf{10}. After the killing, the hunter is proud of his accomplishment and again, the Thought Moments are rooted in \textbf{Attachment}.

We may think the cause of \textbf{Aversion} is outside of us; caused by a person, a situation, etc., but in reality, the cause of \textbf{Aversion} is internal to us. Eleanor Roosevelt\footnote{\url{http://en.wikipedia.org/wiki/Eleanor_Roosevelt}} said, “No one can make you feel inferior without your consent.” \footnote{\url{http://www.goodreads.com/author/quotes/44566.Eleanor_Roosevelt}} This applies to all forms of \textbf{Aversion}, not just feeling inferior.

The Buddha was once insulted by a Brahmin.\footnote{SN 7.2: \url{http://www.accesstoinsight.org/tipitaka/sn/sn07/sn07.002.than.html}} The Buddha replied, “If you offer food to a guest and the guest refuses it, does the food not then belong to you? Similarly, I do not accept the insults that you have offered, so they are all yours.” Hatred is never conquered by hatred, only by loving-kindness.\footnote{Dhammapada verse 5: \url{http://www.tipitaka.net/tipitaka/dhp/verseload.php?verse=005}}

Here is an analogy about dealing with anger. Imagine that we want to stop the steam coming from a pot of boiling water on the stove. The short-term quick fix, the one that we are immediately drawn to, is to put a lid on the pot to keep the steam from escaping while pressure builds up inside the pot. The person who sees things as they truly are, looks at the big picture and understands that it is the nature of steam to arise as long as there is a source of heat. This person stops the steam by turning off the stove. Imagine that we are angry and getting ready to blow off steam. Short-term quick fix techniques such as “firm resolution,” “considering the Buddha’s example” and “considering the harmful effects of anger” deal with the symptoms. It is impossible to overcome anger using a strategy based on \textbf{Aversion} to the current situation. Only beautiful Mental Factors such as \textbf{Mindfulness}, mettā and \textbf{Understanding} deal with the root cause and turn off anger at the source.

Over the years, I have been asked many times, “How do I get rid of anger?” but nobody has ever asked “How do I stop enjoying my morning coffee?” To quote Rudyard Kipling,\footnote{\url{http://en.wikipedia.org/wiki/Rudyard_Kipling}} we need to “meet with triumph and disaster and treat those two impostors just the same.” \footnote{\url{http://www.poetryfoundation.org/poem/175772}} The Buddha identified eight worldly conditions as gain and loss, status and disgrace, praise and blame, pleasure and pain.\footnote{AN 8.6: \url{http://www.accesstoinsight.org/tipitaka/an/an08/an08.006.than.html}} The saint sees these eight worldly conditions as they truly are; impermanent. But when facing these eight worldly conditions, you and I get caught up in \textbf{Attachment} and \textbf{Aversion}.

\subsubsection*{Delusion-rooted}

All 12 of the Thought Moments in the Danger Zone include the root of \textbf{Delusion}. This is why I mentioned earlier that the entire Danger Zone was engulfed in a dense fog causing mental blindness. Reminds me of a joke... if ignorance is bliss, why aren’t more people happy?

\textbf{Attachment}-rooted Thought Moments have mental blindness working in the background, covering up the nature of the object. If the object is seen as impermanent, how can there be \textbf{Attachment}? \textbf{Aversion}-rooted Thought Moments also have mental blindness working in the background, covering up the nature of the object. If the object is seen as naturally arising because of conditions, how can there be \textbf{Aversion}? \textbf{Attachment} and \textbf{Aversion} cannot arise without \textbf{Delusion} doing its job in the background. \textbf{Delusion} is like a magician who can fool you. However, once you know how the magician performs his tricks, he can no longer deceive and captivate you; you are no longer fooled.

\textbf{Delusion} is also like the director of a film. You never see the director on the screen, but the director influences everything from the background. The director controls the script, actors, lighting, camera angle and the music of the film. Everything is coordinated by the director to manipulate the audience’s emotions; \textbf{Attachment} and \textbf{Aversion}. On the other hand, if you were on the set while they were making the movie, it would be obvious what was really happening.

The head of an acting troupe once asked the Buddha about the rebirth destination for an actor.\footnote{SN 42.2: \url{http://www.accesstoinsight.org/tipitaka/sn/sn42/sn42.002.than.html}} The Buddha replied that an actor could be reborn in hell or as an animal because they cause \textbf{Attachment}, \textbf{Aversion} and  \textbf{Delusion} in other people.\footnote{From the context of this Sutta, it is clear that the Buddha was referring to actors whose focus was entertainment; actors whose intention is to educate the audience regarding the human condition may not be included in this category.} 

\begin{figure}[h]
\centering
\includegraphics[width=1.0\linewidth]{./Diagrams/Actors}
\caption{The Buddha criticized actors for causing \textbf{Attachment}, \textbf{Aversion} and \textbf{Delusion} in other people. Many salespeople and politicians also cause \textbf{Attachment}, \textbf{Aversion} and \textbf{Delusion} in other people.}
\label{fig:Actors}
\end{figure}

Thought Moments \textbf{1}--\textbf{8} have both \textbf{Delusion} and \textbf{Attachment}. Thought Moments \textbf{9} and \textbf{10} have both \textbf{Delusion} and \textbf{Aversion}. Thought Moments \textbf{11} and \textbf{12} have \textbf{Delusion}, but no \textbf{Attachment} and no \textbf{Aversion}.

Thought Moment \textbf{11} is associated with \textbf{Doubt}. \textbf{Doubt} is not the same as not being sure about a person’s name. \textbf{Doubt} refers to the uncertainty in the benefit of generosity, uncertainty about the benefit of morality or uncertainty about the benefit of spiritual development. At the end of the last talk, I said that the Theravāda Suttas we have today contain the teachings of the Buddha, but they are \textbf{not} the literal, verbatim “word of the Buddha.” This qualifies under the general English language usage of the word “doubt,” but does not imply Thought Moment \textbf{11}.

\textbf{Doubt} cannot be reduced just by thinking; direct experience is the way to reduce \textbf{Doubt}. Here is a metaphor to show how \textbf{Doubt} is uprooted by the Sotāpanna’s experience of \textit{Nibbāna}. Imagine that you are on the near side of a stream and want to get to the far side.\footnote{The near side of the stream is Self-view; the Sotāpanna has eradicated Self-view. See SN 35.197: \url{http://www.accesstoinsight.org/tipitaka/sn/sn35/sn35.197.than.html}} You start by convincing yourself saying, “I can jump over this stream.” This is blind faith; there is still a lot of \textbf{Doubt} in the mind. You study books on stream jumping, you listen to talks on how to jump across streams, you read about previous stream-jumpers and now when you say, “I can jump over this stream,” there is less \textbf{Doubt}. On your side of the stream you practice jumping and become incredibly skilled, an Olympic-level jumper. Now when you say, “I can jump over this stream,” there is much less \textbf{Doubt}, but \textbf{Doubt} remains because you have not yet jumped over this stream. You jump over it and are now on the far side. Having jumped over it, there are no more conditions for \textbf{Doubt} to arise. There is not even the most subtle \textbf{Doubt} remaining. \textbf{Doubt} has been truly uprooted. Before this point, \textbf{Doubt} was increasingly suppressed, but now and forever in the future, it has been uprooted.

The handout indicates that Thought Moment \textbf{12} is associated with \textbf{Restlessness}. Actually, \textbf{Restlessness} arises in all of the unwholesome Thought Moments, but in Thought Moments \textbf{1}--\textbf{11}, \textbf{Restlessness} is more in the background whereas in Thought Moment \textbf{12}, it is dominant.

\textbf{Restlessness} is the opposite of steadiness or calmness; it is confusion, unsteadiness, agitation, mental distraction or mental excitement. A restless mind is not interested in the object and treats the object superficially.

\textbf{Restlessness} arises very frequently, but most people are unaware of it. When a person first tries to meditate, they are often surprised at the \textbf{Restlessness} in the mind. It seems that the mind never wants to penetrate the object; that the mind is wild and untamed. To help meditators deal with \textbf{Restlessness}, the Buddha delivered a Sutta titled “The removal of distracting thoughts.” \footnote{MN 20: \url{http://www.accesstoinsight.org/tipitaka/mn/mn.020.soma.html}} The good news is that the \textbf{Volition} associated with Thought Moment \textbf{12} is very weak, so the unwholesome kamma created is not weighty.

\subsection*{Thought Moments 31--38 (Faultless Zone)}

\begin{figure}[h]
\centering
\includegraphics[width=0.46\linewidth]{./Diagrams/Faultless}
\caption{Thought Moments \textbf{31}--\textbf{38}: Faultless Zone}
\label{fig:Faultless}
\end{figure}

Let’s now jump to Thought Moments that create wholesome kamma, Thought Moments \textbf{31}--\textbf{38}. Earlier, I called this the Faultless Zone. This group is subdivided based on “prompted and unprompted,” “with and without \textbf{Understanding}” and “pleasant and indifferent \textbf{Feeling}.” We have already discussed prompted and unprompted, and pleasant and indifferent \textbf{Feeling}, so let’s now discuss \textbf{Understanding}.

In Pāḷi, the word for \textbf{Understanding} is \textit{Paññā}, the Abhidhamma considers \textit{Paññā} to be the foundation of wisdom, investigation, insight, right view and \textit{vipassanā}.\footnote{One Mental Factor listed in the Abhidhamma is often equivalent to multiple terms from the Suttas.} Thought Moments \textbf{31}--\textbf{38} will create wholesome kamma, but Thought Moments associated with \textbf{Understanding}, Thought Moments \textbf{31}, \textbf{32}, \textbf{35} and \textbf{36} will contribute to spiritual development.\footnote{Faultless Thought Moments without \textit{Paññā} may also contribute to spiritual development, depending on how one defines “spiritual development.”}

Just to be clear, \textbf{Understanding} or wisdom does not imply intelligence. There is a story in the Commentary of a monk who was unable to memorize a single verse after trying for four months, yet was able to become an Arahat.\footnote{\url{http://www.tipitaka.net/tipitaka/dhp/verseload.php?verse=025}}

Morality and thoughts of generosity always create wholesome kamma, and when associated with \textbf{Understanding}, this can also contribute to spiritual development. All forms of mental cultivation contribute to spiritual development. This includes meditation, studying the Dhamma, teaching the Dhamma and straightening of views.

\begin{figure}[h]
\centering
\includegraphics[width=0.6\linewidth]{./Diagrams/Sila}
\caption{Latent tendencies give rise to thoughts which are manifested as actions and speech. Sīla is the restraint of actions and speech, Samadhi calms the mind and Paññā exposes the latent tendencies.}
\label{fig:Sila}
\end{figure}

So how do we apply \textbf{Understanding} in our daily life?\footnote{Taken from Visuddhimagga I.13 (see footnote 2).} Bad speech and bad action come from bad thoughts. Bad thoughts arise from latent tendencies. Working backwards, we start by addressing bad speech and bad actions through precepts, morality or virtue; in Pāḷi, this is called \textit{Sīla}. Bad thoughts are addressed through having a composed or concentrated mind; in Pāḷi, this is called \textit{Samādhi}. Latent tendencies are addressed through \textbf{Understanding} or \textit{Paññā}.

The Noble Eightfold Path can also be analyzed in terms of \textit{Sīla}, \textit{Samādhi} and \textit{Paññā}. \textit{Sīla} is Right Speech, Right Action and Right Livelihood. \textit{Samādhi} is Right Effort, Right \textbf{Mindfulness} and Right Concentration. \textit{Paññā} is Right View and Right Thought.

\textit{Sīla}, \textit{Samādhi} and \textit{Paññā} are progressive steps. Start with \textit{Sīla}, and that will support \textit{Samādhi} which is a condition for \textit{Paññā}. \textit{Paññā} deepens the quality of \textit{Sīla} leading to deeper \textit{Samādhi} and clearer \textit{Paññā}, and the cycle continues.

The Commentary gives an analogy to explain the difference between perceiving, thinking and \textbf{Understanding} to emphasize that \textbf{Understanding} is not superficial. A baby, a child and a money-changer all see a coin. The baby perceives a shiny, round object. The child thinks, “This is money” and thinks about what money can buy, how money can be used. The money-changer understands the qualities of the coin and can detect a counterfeit coin. The baby perceives, the child thinks and the money-changer \textbf{Understands}.

The Buddha talked about the blind men who each touch different parts of an elephant and then argue over “what is an elephant.” \footnote{Ud 6.4: \url{http://www.accesstoinsight.org/tipitaka/kn/ud/ud.6.04.than.html}} The blind men’s different perceptions led to incomplete thinking and arguments. \textbf{Understanding} would be viewing the big picture like the sighted person who can see the entire elephant. Reminds me of a joke... a group of blind elephants each touched different parts of a human and they all agreed that a human was flat.

Earlier, I mentioned how \textbf{Attachment} to \textbf{Wrong view} comes from distorted perception and distorted thinking. The analogies of the snake, the coin and the elephant all remind us that we all are subject to distorted and incomplete perceptions, distorted and incomplete thinking. Assumptions can lead to \textbf{Wrong view} but investigation can lead to \textbf{Understanding}.

Let’s consider another example to identify the associated Thought Moment. A mother sees a monk on alms round through the window. Immediately, she collects some food and takes her young child outside. They both offer food to the monk. At the time of offering, the Thought Moment of the mother and the Thought Moment of the young child are both wholesome; the Thought Moments are in the Faultless Zone. The Thought Moment of the mother is unprompted, spontaneous, while the Thought Moment of the child is prompted and depended on the inducement of the mother. The Thought Moment of the mother is with \textbf{Understanding}; she sees this as part of her spiritual development, and the child’s Thought Moment is without \textbf{Understanding}, the child is too young to comprehend. The Thought Moment of the mother is accompanied by pleasant \textbf{Feeling} while the child’s Thought Moment may be accompanied by indifferent \textbf{Feeling}. Mother and child perform exactly the same action, but the mother’s kamma will be much weightier because her \textbf{Volition} was much stronger. The mother’s Thought Moment will be \textbf{31} and the child’s will be \textbf{38}.

\subsection*{Thought Moments 13--30 (Sensing Zone)}

\begin{figure}[h]
\centering
\includegraphics[width=0.46\linewidth]{./Diagrams/Sensing}
\caption{Thought Moments \textbf{13}--\textbf{30}: Sensing Zone}
\label{fig:Sensing}
\end{figure}

We have covered the Danger Zone, Thought Moments \textbf{1}--\textbf{12} and the Faultless Zone, Thought Moments \textbf{31}--\textbf{38}. Now let’s look at the Sensing Zone, Thought Moments \textbf{13}--\textbf{30}. These are involved in the sensing and processing of sense data and ideas. In a later talk, I will go into the details of each of these Thought Moments, but at this point I will give only a brief overview of some of them.

Thought Moments \textbf{13}--\textbf{17} and Thought Moments \textbf{20}--\textbf{24} are the sense-consciousness Thought Moments corresponding to the five physical senses. The Thought Moment that captures an image at the back of the retina can be either Thought Moment \textbf{13} or Thought Moment \textbf{20}; both are eye-consciousness Thought Moments. The only difference between them is that Thought Moment \textbf{13} is the result of some past unwholesome kamma, while Thought Moment \textbf{20} is the result of some past wholesome kamma.

Once the sense-consciousness Thought Moment has performed its function, the sense data is processed. Thought Moments \textbf{18} and \textbf{19} will process sense data that arose because of past unwholesome kamma. Thought Moments \textbf{25} and either Thought Moment \textbf{26} or Thought Moment \textbf{27} will process sense data that arose because of past wholesome kamma.

Once sense data has been captured and processed, a decision needs to be made as to how to proceed. That decision is the function of Thought Moment \textbf{29}. If the object is an idea, then Thought Moments \textbf{13}--\textbf{27} are bypassed, and the mind jumps directly to Thought Moment \textbf{29} to make a decision. You will notice that Thought Moment \textbf{29} is unrelated to kamma. The mind does not make decisions because of kamma.

\subsection*{Room analogy}

\begin{figure}[h]
\centering
\includegraphics[width=0.4\linewidth]{./Diagrams/Room}
\caption{Decision directs the flow of the mind from the input/trigger to one of the outputs.}
\label{fig:Room}
\end{figure}

Let me give you an analogy to explain the interaction between the Danger Zone, the Faultless Zone and the Sensing Zone. There is a diagram on the handout titled “Room Analogy.” There is one entrance to the room on the south wall; three exits from the room on the west wall, and one exit from the room on the east wall. The mind comes into the room from the south, and leaves either through one of the three exits on the left or through the exit on the right.

In this analogy, the input or trigger at the entrance is how a stimulus such as a \textbf{Visible-form}, \textbf{Sound}, idea, etc. enters the mind. The exits are the mind’s reaction to the stimulus; the result of the decision made by Thought Moment \textbf{29}. The three exits to the left are based on \textbf{Attachment}, \textbf{Aversion} and \textbf{Delusion}, and lead to the Danger Zone. The exit to the right leads to the Faultless Zone.

In the previous talk, I mentioned that the mind makes decisions naturally. I talked about the experiments that show how the mind decides first, and then later the idea of a Self who decides is created. I mentioned that a concept of a Self is added by the mind as a justification or rationalization of the decision making process. Actually, it is the decisions made by Thought Moment \textbf{29} that determines how much time is spent in the Faultless Zone and how much time is spent in the Danger Zone.

So how does Thought Moment \textbf{29} decide? The answer is very simple: accumulation of habits, not kamma. Accumulated good habits cause Thought Moment \textbf{29} to open the door to the Faultless Zone. Accumulated bad habits cause Thought Moment \textbf{29} to open one of the doors to the Danger Zone. So how can accumulated habits be changed or reinforced? The answer is very simple: training.

Simple is not the same as easy. The five precepts, the rules of training, are simple to understand, but keeping the precepts is sometimes difficult. I try to recite the five precepts at least once a day. When I recite them, I reflect deeply upon them. Repetitive actions done with strong \textbf{Volition} trains the mind and reinforces habits.

\subsection*{Thought Moments 55--81}

\begin{figure}[h]
\centering
\includegraphics[width=0.5\linewidth]{./Diagrams/55to81}
\caption{Thought Moments \textbf{55}--\textbf{81}: Fine Material Sphere and Immaterial Sphere}
\label{fig:55to81}
\end{figure}

Let’s move to Thought Moments \textbf{55}--\textbf{69}, which relate to the jhānas. As the focus of these talks is applying Abhidhamma to daily life, I will not spend a lot of time discussing jhāna practice. Jhāna practice pre-dates the Buddha; the Buddha talked about practising the jhānas while still a Bodhisatta, before he was enlightened.\footnote{MN 26: \url{http://www.accesstoinsight.org/tipitaka/mn/mn.026.than.html}} Jhāna practice is part of the Buddha’s teaching; the definition of “Right Concentration” as part of the Noble Eightfold Path is the practice of the four jhānas.\footnote{SN 45.8: \url{http://www.accesstoinsight.org/tipitaka/sn/sn45/sn45.008.than.html}}

A slight paraphrase of the standard description of jhāna found in the Suttas reads as follows: “Before entering the first jhāna, one must temporarily subdue the senses and temporarily subdue the unwholesome Mental Factors of sense desire, ill will, \textbf{Sloth} and \textbf{Torpor}, \textbf{Restlessness} and worry, and \textbf{Doubt} which block the first jhāna from arising. The first jhāna arises with the Mental Factors of \textbf{Zest}, pleasant \textbf{Feeling}, thought, examination and \textbf{One-pointedness}. When thought and examination fall away, one enters the second jhāna. When \textbf{Zest} falls away, one enters the third jhāna. When pleasant \textbf{Feeling} changes to indifferent \textbf{Feeling}, one enters the fourth jhāna.”

I am sometimes asked, “The Suttas always talk about four jhānas, so why does the Abhidhamma list five?” The Suttas describe the first jhāna as having \textbf{Zest}, pleasant \textbf{Feeling}, thought, examination and \textbf{One-pointedness}, and the transition from the first jhāna to the second jhāna drops both thought and examination. But there is a Sutta that describes concentration without thought but with examination, kind of part way between the first and second jhāna.\footnote{DN 33: \url{http://suttacentral.net/en/dn33}} The Abhidhamma defines this “part way between” jhāna as the second jhāna and therefore the second jhāna according to the Suttas becomes the third jhāna according to the Abhidhamma. With this additional jhāna according to the Abhidhamma, the third jhāna according to the Suttas is the fourth jhāna according to the Abhidhamma and the fourth jhāna is the fifth jhāna according to the Abhidhamma.

\begin{figure}[h]
\centering
\includegraphics[width=0.4\linewidth]{./Diagrams/Jhana}
\caption{Jhānas in the Abhidhamma and in the Suttas.}
\label{fig:Jhana}
\end{figure}

Actually, the Suttas describe the second jhāna, without thought and without examination, as “Noble Silence” because there is no more mental chatter.\footnote{SN 21.1: \url{http://www.accesstoinsight.org/tipitaka/sn/sn21/sn21.001.than.html}} Remember this when you see a “Noble Silence” sign in a meditation hall: perhaps they are asking you to achieve the second jhāna!

Thought Moments \textbf{55}--\textbf{59} arise when experiencing the first to fifth jhāna. These jhānas are analogous to being in an aeroplane; a special kind of environment detached from what is happening on the ground.

Using the fifth jhāna as a foundation, the meditator turns to increasingly subtle objects: infinite space, infinite consciousness, nothingness and “neither-perception-nor-non-perception.” These are Thought Moments \textbf{70}--\textbf{73}. These jhānas are analogous to being in a spaceship; an even more specialized environment, even more detached from what is happening on the ground.

\subsection*{Thought Moments 82--89\footnote{Handout 2 lists 89 Thought Moments, including 81 mundane Thought Moments and 8 supramundane Thought Moments. Sometimes the 8 supramundane Thought Moments are each subdivided into 5 according to the accompanying jhāna. When this is done, there are 81 mundane Thought Moments and 40 supramundane Thought Moments for a total of 121.}}

\begin{figure}[h]
\centering
\includegraphics[width=0.19\linewidth]{./Diagrams/82to89}
\caption{Thought Moments \textbf{82}--\textbf{89}: Supramundane}
\label{fig:82to89}
\end{figure}

Thought Moments \textbf{82}--\textbf{89} are the “world transcending” or Supramundane Thought Moments. As I mentioned earlier, these are the only Thought Moments that can take \textit{Nibbāna} as object and are reserved for saints and for the Buddha.

A fetter is a chain around the ankles of prisoners. There are 10 fetters that bind beings to \textit{saṃsāra}.\footnote{AN 10.13: \url{http://www.accesstoinsight.org/tipitaka/an/an10/an10.013.than.html}} The first five, the ones that bind beings to the sensuous realm, are personality-belief, \textbf{Doubt}, \textbf{Attachment} to rules and rituals, sensuous craving and ill will.

\begin{figure}[h]
\centering
\includegraphics[width=1.0\linewidth]{./Diagrams/Magga}
\caption{Thought Moment \textbf{82} is the transition to the purified mind of a Sotāpanna.}
\label{fig:Magga}
\end{figure}

The arising of Thought Moment \textbf{82} in the mental stream of an individual represents a “change of lineage;” that being is now a Sotāpanna.\footnote{According to the Suttas (AN 3.86: \url{http://www.accesstoinsight.org/tipitaka/an/an03/an03.086.than.html}), a Sotāpanna will be reborn a maximum of seven times (never in a Woeful Plane) before becoming and Arahat.} This is the first experience of \textit{Nibbāna}, and this experience impacts all future Thought Moments, including Thought Moments in future lives. After this experience, there are no more conditions to support the arising of the first three fetters. In other words, it is impossible for personality-belief, \textbf{Doubt} or \textbf{Attachment} to rules and rituals to arise in a Sotāpanna’s mind.

Earlier, I gave the analogy of jumping over the stream. Having already jumped over the stream, there are no conditions for \textbf{Doubt} to arise such as “can I do it or not?” Because of this first experience of \textit{Nibbāna}, there are no more conditions for \textbf{Doubt} to arise in the mind of a Sotāpanna.

I also talked earlier about the man with a latent fear of snakes, and how distortion of perception leads to distortion of thinking, distortion of thinking leads to distortion of views and distortion of view leads to unwholesome speech and action. The Sotāpanna’s first experience of \textit{Nibbāna} convinces him that everything is just mind and matter (\textit{nāmarūpa}), and that there is no underlying Self. Distortions of perception and distortions of thinking may still arise in a Sotāpanna, but there are no more conditions for distortion of views and therefore, no speech or action that could lead to rebirth in a woeful state such as hell or as an animal.

The third fetter uprooted by the Sotāpanna’s first experience of \textit{Nibbāna} is \textbf{Attachment} to rules and rituals. The Sotāpanna is convinced that speech and action such as chanting, or bowing in front of the Buddha statue, are not required to experience \textit{Nibbāna}. The Sotāpanna is convinced that development of the mind, not speech or action, allows one to experience \textit{Nibbāna}. Just to be clear, I am not suggesting that chanting or bowing is bad; I am saying that \textbf{attachment} to precepts or worship does not lead to spiritual development.

The remaining two fetters that bind us to the sensuous realm, sensuous craving and ill will, are significantly weakened by the Sakadāgāmī and uprooted by the Anāgāmī. When a Sotāpanna experiences Thought Moment \textbf{83}, he is “promoted” to the level of Sakadāgāmī, and when an Sakadāgāmī experiences Thought Moment \textbf{84}, he is “promoted” to the level of Anāgāmī.

When an Anāgāmī experiences Thought Moment \textbf{85}, he is “promoted” to the level of Arahat, and five more fetters are uprooted. The five fetters are: craving for fine-material existence, craving for immaterial existence, \textbf{Conceit}, \textbf{Restlessness} and \textbf{Delusion}.

Sometimes, I am asked, “If a Sotāpanna has eliminated personality belief, how can he still have \textbf{Conceit}?” Though a Sotāpanna sees beings as just \textit{nāmarūpa}, he may still compare this set of \textit{nāmarūpa} with that set of \textit{nāmarūpa}. \textbf{Conceit} is any form of comparison.

Once Thought Moment \textbf{82} has arisen, there is no condition to support the arising of Thought Moments \textbf{1}, \textbf{2}, \textbf{5}, \textbf{6} and \textbf{11}. Once Thought Moment \textbf{84} has arisen, there is no condition to support the arising of Thought Moments \textbf{9} and \textbf{10}. An Arahat cannot have any unwholesome Thought Moments and does not create new kamma, so once Thought Moment \textbf{85} has arisen, there is no condition to support the arising of Thought Moments \textbf{1}--\textbf{12}, \textbf{31}--\textbf{38}, \textbf{55}--\textbf{59} or \textbf{70}--\textbf{73}.

A Sotāpanna can experience the “bliss of \textit{Nibbāna}.” Thought Moments lower than 82 cannot take \textit{Nibbāna} as an object, and Thought Moment \textbf{82} arises only once, when the ordinary being becomes a Sotāpanna. Thought Moment \textbf{86} is the Sotāpanna fruit; it takes \textit{Nibbāna} as an object. So when a Sotāpanna experiences the “bliss of \textit{Nibbāna},” Thought Moment \textbf{86} is arising. When a Sakadāgāmī experiences the “bliss of \textit{Nibbāna},” Thought Moment \textbf{87} is arising. When an Anāgāmī experiences the “bliss of \textit{Nibbāna},” Thought Moment \textbf{88} is arising. And finally, when an Arahat experiences the “bliss of \textit{Nibbāna},” Thought Moment \textbf{89} is arising. One of the blessings sometimes used by monks is, “May this be a condition leading to path and fruit knowledge.” What they are saying is, “May this be a condition leading to sainthood.”

\subsection*{Categories of Thought Moments}

At the beginning of this talk, I gave a brief description of 13 groups of Thought Moments in the handout. This grouping is only one way of categorizing the Thought Moments.

A second way of categorizing them is according to kamma. The Thought Moments in the top row create new kamma. The Thought Moments in the middle row are the result of past kamma, and the Thought Moments in the bottom row are unrelated to kamma.

A third way of categorizing Thought Moments is according to roots. Thought Moments \textbf{1}--\textbf{12} have unwholesome roots, Thought Moments \textbf{13}--\textbf{30} have no roots, and the remaining Thought Moments have beautiful roots. The beautiful roots are \textbf{Non-attachment}, \textbf{Non-aversion} and \textbf{Understanding}.

A fourth way of categorizing Thought Moments is according to sphere. We will discuss this during a later talk, but you can see from the handout that Thought Moments \textbf{1}--\textbf{54} are in the Sense Sphere category, Thought Moments \textbf{55}--\textbf{69} are in the Fine Material sphere category, Thought Moments \textbf{70}--\textbf{81} are in the Immaterial Sphere category, and Thought Moments \textbf{82}--\textbf{89} are in the Supramundane category.

A fifth way of categorizing Thought Moments is according to \textbf{Feeling}. The accompanying \textbf{Feeling} is shown in the handout for Thought Moments \textbf{1}--\textbf{54}, but the remaining Thought Moments also have accompanying \textbf{Feeling}.

The first book of the Abhidhamma lists more than 40 ways of classifying Thought Moments according to the Suttas, and more than 120 ways of classifying Thought Moments according to the Abhidhamma. The Abhidhamma is very thorough!

\subsection*{Linkage to \textit{Satipaṭṭhāna} Sutta}

\subsubsection*{Contemplation of \textbf{Feeling}}

Let’s now discuss the contemplation of \textbf{Feeling} from the \textit{Satipaṭṭhāna} Sutta. In paragraph 32, when the Sutta says “contemplating \textbf{Feelings} in \textbf{Feelings},” it means the meditator should repeatedly observe \textbf{Feeling}, but limit the involvement to only observing the \textbf{Feeling}.\footnote{The Pāḷi word for “contemplate” is \textit{anupassi}; The prefix \textit{anu} gives emphasis and \textit{passati} is the Pāḷi verb, “to see.” The Visuddhimagga XXI.14 (see footnote 2) explains this as “sees again and again in various modes.”} In other words, when observing pleasant \textbf{Feeling}, do not slip into \textbf{Attachment}; when observing unpleasant \textbf{Feeling}, do not slip into \textbf{Aversion}; when observing neutral \textbf{Feeling}, do not slip into \textbf{Restlessness}. Observing \textbf{Attachment}, \textbf{Aversion} and \textbf{Restlessness} is part of the practice of contemplation of consciousness, but this is not the practice of contemplation of \textbf{Feeling}.

Paragraph 33 asks the meditator to observe pleasant, unpleasant and neutral \textbf{Feeling} in general, and then asks the meditator to observe pleasant, unpleasant and neutral \textbf{Feeling} according to their source. The two sources of \textbf{Feelings} are “worldly” and “spiritual.” Worldly \textbf{Feelings} arise from \textbf{Contact} with the six sense objects. Spiritual \textbf{Feelings} are related to spiritual development such as generosity, morality, mental cultivation and jhāna.

Spiritual pleasant \textbf{Feeling} arises with Thought Moments \textbf{31} and \textbf{32} as well as the jhāna Thought Moments, accompanied by pleasant \textbf{Feeling}. Spiritual neutral \textbf{Feeling} arises with Thought Moments \textbf{35} and \textbf{36} as well as the jhāna Thought Moments, accompanied by neutral \textbf{Feeling}. Spiritual unpleasant \textbf{Feeling} can arise at different stages of insight. For example, the Commentary mentions the contemplation of pain to abandon the perception of pleasure; during this practice, spiritual unpleasant \textbf{Feeling} will arise.\footnote{Visuddhimagga XX.90 (see footnote 2).}

\subsubsection*{Contemplation of Consciousness}

\begin{figure}[h]
\centering
\includegraphics[width=0.7\linewidth]{./Diagrams/Aversion}
\caption{The mind with \textbf{Aversion} often spins out of control, fixing itself to ideas which lead to more \textbf{Aversion}. Contemplation of Consciousness is \textbf{Mindfulness} of a recent Thought Moment (which may have had \textbf{Aversion}). A Thought Moment with \textbf{Mindfulness} is wholesome. In this way, an unwholesome Thought Moment can be a condition for a wholesome Thought Moment.}
\label{fig:Aversion}
\end{figure}

Moving on to the contemplation of consciousness, please look at paragraph 36. When the Sutta says, “know the consciousness with lust, as with lust,” we are being asked to recognize Thought Moments \textbf{1}--\textbf{8}. We can recognize them because they are sticky. When the Sutta says, “know the consciousness without lust, as without lust,” we are being asked to recognize Thought Moments other than \textbf{1}--\textbf{8}. The consciousness with hate are Thought Moments \textbf{9} and \textbf{10}. We can recognize them because they burn. The consciousness with ignorance are Thought Moments \textbf{1}--\textbf{12}. We can recognize them because they are foggy.

The shrunken state refers to Thought Moments that are prompted, not spontaneous, or as footnote 13 at the end of the Sutta says, “a rigid and indolent state of mind.” The distracted state refers to Thought Moments with \textbf{Restlessness}. As mentioned earlier, \textbf{Restlessness} arises in all unwholesome Thought Moments, but in Thought Moments \textbf{1}--\textbf{11}, \textbf{Restlessness} is more in the background whereas in Thought Moment \textbf{12}, restless is dominant.

The final four classifications of consciousness, “the developed state,” “the state with nothing higher,” “the concentrated state,” and “the freed state,” all refer to monitoring the mind that is in advanced states of meditative development such as the jhānas.

\subsection*{Summary of Key Points}

Here is a summary of key points regarding Consciousness:

\begin{itemize}

\item A Thought Moment consists of consciousness and a collection of Mental Factors; within the Thought Moment, consciousness has the function of awareness, and each of the Mental Factors performs its own individual function.

\item Handout 2 shows the map of the mind, consisting of 89 Thought Moments:

\begin{itemize}

\item Thought Moments \textbf{1}--\textbf{12} (“Danger Zone”) create new unwholesome kamma.

\begin{itemize}

\item Thought Moments \textbf{1}--\textbf{8} are \textbf{Attachment}-rooted (sticky quicksand) and are classified according to:

\begin{itemize}

\item \textbf{Feeling} (pleasant/indifferent).

\item Associated with \textbf{Attachment} to \textbf{Wrong view}/not associated with \textbf{Attachment} to \textbf{Wrong view}.

\item Unprompted (spontaneous)/prompted (induced).

\end{itemize}

\item Thought Moments \textbf{9}--\textbf{10} are \textbf{Aversion}-rooted (hot and painful) and are accompanied by unpleasant \textbf{Feeling}.

\item Thought Moments \textbf{11}--\textbf{12} are \textbf{Delusion}-rooted (thick fog); Thought Moment \textbf{11} is with \textbf{Doubt}, Thought Moment \textbf{12} is with \textbf{Restlessness}.

\end{itemize}

\item Thought Moments \textbf{13}--\textbf{30} (“Sensing Zone”) process sense data.

\item Thought Moments \textbf{31}--\textbf{54} (“Faultless Zone”) includes Thought Moments \textbf{31}--\textbf{38}, which create new wholesome kamma.

\begin{itemize}

\item These Thought Moments are classified according to:

\begin{itemize}

\item \textbf{Feeling} (pleasant/indifferent).

\item Associated with \textbf{Understanding}/not associated with \textbf{Understanding}.

\item Unprompted (spontaneous)/prompted (induced).

\end{itemize}

\end{itemize}

\item Thought Moments \textbf{55}--\textbf{81} are related to jhāna meditative states.

\item Thought Moments \textbf{82}--\textbf{89} include the attaining of the four degrees of Sainthood and enjoying the bliss of \textit{Nibbāna}.

\end{itemize}

\item The Room Analogy shows the mind “deciding” between the “Danger Zone” and the “Faultless Zone” when triggered by the “Sensing Zone.”

\item Handout 2 helps us to \textbf{\textit{R}}ecognize the current Thought Moment; \textbf{\textit{R}}ecognize is the first step in the \textbf{\textit{R}} \textbf{\textit{A}} \textbf{\textit{D}} \textbf{\textit{I}} \textbf{\textit{CA}} \textbf{\textit{L}} process (\textbf{\textit{R}}ecognize, \textbf{\textit{A}}ccept, \textbf{\textit{D}}epersonalize, \textbf{\textit{I}}nvestigate, \textbf{\textit{C}}ontemplate \textbf{\textit{A}}\textit{nicca}/\textbf{\textit{C}}ontemplate \textbf{\textit{A}}\textit{nattā}, \textbf{\textit{L}}et go).

\end{itemize}

Finally, in my opinion, the most important thing to remember about consciousness and Thought Moments is that we should know if the mind is in the Danger Zone or in the Faultless Zone. If the mind is in the Danger Zone, reflect upon the disadvantages of such thinking. If the mind is in the Faultless Zone, just be passively aware.

\begin{center}
\textbf{\textit{This concludes the third talk.}} \\
\end{center}

\newpage

\subsection*{Questions \& Answers}
\color {red}

\question{What is the duration of a Thought Moment?}

According to the commentaries, more than a billion Thought Moments occur in the time occupied by a flash of lightning.\footnote{See entry for \textit{citta-kkhaṇa} in “Buddhist Dictionary” (see Footnote 2).} The Buddha said that he could not come up with a simile to describe how fast the mind changed.\footnote{AN 1.48: \url{http://www.accesstoinsight.org/tipitaka/an/an01/an01.048.than.html}}

We may believe that our senses are working in parallel, that we see, hear and think at the same time, but according to the Abhidhamma, these processes happen in very fast succession. We have heard stories of people whose life has flashed before their eyes in a few seconds. This suggests that under certain circumstances, the mind can be aware of a lot of precise details in a very short period of time. When observing the mind, experienced meditators can observe a great level of detail.

\question{Is the Abhidhamma useful for  “socially-engaged Buddhists”?}

The Abhidhamma helps us recognize when the mind is in the Danger Zone, under the influence of the three unwholesome roots of \textbf{Attachment}, \textbf{Aversion} and \textbf{Delusion}. Unfortunately, these three roots have become institutionalized in today's society. The root of Attachment is nourished by our culture of consumerism. The root of Aversion is nourished by 

\color {black}

\question{How does the Abhidhamma explain the concepts of “subconscious” and the “unconscious mind?”}

The Wikipedia article\footnote{\url{http://en.wikipedia.org/wiki/Subconscious}} on “subconscious” explains that subconscious is the part of consciousness that is not currently in focal awareness, and that since there is a limit to what can be held in conscious focal awareness, an alternative storehouse of one’s knowledge and prior experience is needed.\footnote{Quote by Carl Jung (\url{http://en.wikipedia.org/wiki/Carl_Jung}), who wrote about the subconscious.} According to the Abhidhamma, there is only one Thought Moment at a time, and this Thought Moment includes consciousness or awareness. Prior experience influences the Thought Moment through natural decisive support condition, which will be explained in a later talk.

The Wikipedia article\footnote{\url{ http://en.wikipedia.org/wiki/Unconscious_mind}} on the “unconscious mind” defines it as consisting of the processes in the mind that occur automatically and are not available to introspection, including memory and motivation; though these processes exist well under the surface of conscious awareness, they are theorized to exert an impact on behaviour.\footnote{Sigmund Freud (\url{ http://en.wikipedia.org/wiki/Sigmund_Freud}) popularized this concept.} This is another way of describing the effect of natural decisive support condition.

\question{Would it be correct to say that on Handout 2, the top row labelled “Create New Kamma” is linked to the middle row labelled “Result of Past Kamma?”}

The Life-continuum Thought Moment \textbf{60} is always the result of Thought Moment \textbf{55} arising in a previous life. In other words, to be reborn in a Brahma Realm, it is necessary to have attained the first jhāna in a previous existence. It is not possible to draw the same kind of connections involving Sense Sphere Thought Moments. For example, it is not possible to say that Thought Moment \textbf{39} is the result of Thought Moment \textbf{31}. As discussed in the talk on Realms of Existence, it is actually much more complicated than this.
